% Part: sets-functions-relations
% Chapter: functions
% Section: basics

\documentclass[../../../include/open-logic-section]{subfiles}

\begin{document}

\olfileid[tr]{sfr}{fun}{bas}
\olsection{Temel Kavramlar}

\begin{explain}
Bir \emph{fonksiyon}, verilen bir kümenin her elemanını verilen başka bir
kümedeki belirli bir elemana gönderen eşlemedir. Örneğin $1$ ekleme işlemi bir
fonksiyon tanımlar: her $n$ sayısı biricik $n+1$ sayısına eşlenir.

Daha genel olarak fonksiyonlar girdi olarak ikililer, üçlüler vb. alıp bir tür
çıktı verebilir. Temel aritmetikten birçok fonksiyonu tanırız. Örneğin toplama
ve çarpma birer fonksiyondur. İki sayı alıp üçüncü bir sayı verirler.

Bu matematiksel, soyut anlamda fonksiyon bir \emph{kara kutu}dur: önemli olan
çıktının hesaplanma yöntemi değil, hangi çıktının hangi girdiye eşlendiğidir.
\end{explain}

\begin{defn}[Fonksiyon]
Bir $f \colon A \to B$ \emph{fonksiyonu}, $A$'nın her elemanını $B$'nin bir
elemanına eşleyen bir eşlemedir.

$A$'ya $f$'nin \emph{tanım kümesi}, $B$'ye ise $f$'nin \emph{değer kümesi}
deriz. $A$'nın elemanlarına $f$'nin girdileri veya \emph{argümanları};
$f$'nin bir $x$ argümanıyla eşlediği $B$ elemanına da $f$'nin $x$
argümanındaki \emph{değeri} denir ve bu değer $f(x)$ ile yazılır.

$f$'nin \emph{görüntü kümesi} $\ran{f}$, değer kümesinin $f$'nin bir
argümandaki değerlerinden oluşan alt kümesidir; $\ran{f} =
\Setabs{f(x)}{x \in A}$.
\end{defn}

\olref{fig:function}'teki şema, fonksiyonları düşünmeyi kolaylaştırabilir.
Soldaki elips fonksiyonun \emph{tanım kümesini}, sağdaki elips fonksiyonun
\emph{değer kümesini} temsil eder; bir ok da tanım kümesindeki bir
\emph{argümandan} değer kümesindeki karşılık gelen \emph{değere} yönelir.

\begin{figure}
  \olasset{assets/diagrams/function.tikz}
  \caption{Bir fonksiyon, bir kümenin her elemanını başka bir kümenin bir
    elemanına eşler. Bir ok, tanım kümesindeki bir argümandan değer
    kümesindeki karşılık gelen değere yönelir.}
  \ollabel{fig:function}
\end{figure}

\begin{ex}
Çarpma, girdi olarak doğal sayı ikililerini alıp çıktı olarak doğal sayılara
eşler; dolayısıyla $\Nat \times \Nat$'tan (tanım kümesinden) $\Nat$'a (değer
kümesine) gider. Her $n \in \Nat$ sayısı $n \times 1$ olduğundan, görüntü
kümesi de $\Nat$'tır.
\end{ex}

\begin{ex}
Çarpma bir fonksiyondur; çünkü her girdiyi, yani her doğal sayı ikilisini,
tek bir çıktıyla eşler: $\times \colon \Nat^2 \to \Nat$. Buna karşılık, bir
$n$ doğal sayısını $x^2=n$ olacak biçimde bir $x$ gerçel sayısına eşlemek
fonksiyonel değildir; çünkü her pozitif tam sayı $n$'nin $\sqrt{n}$ ve
$-\sqrt{n}$ olmak üzere iki karekökü vardır. Yalnızca negatif olmayan esas
karekökü vererek bunu fonksiyonel hâle getirebiliriz:
$\sqrt{\phantom{X}} \colon \Nat \to \Real$.
\end{ex}

\begin{ex}
Bir sınıftaki her öğrenciyi dönem sonu notuyla eşleyen bağıntı bir
fonksiyondur; aynı sınıfta hiçbir öğrenci iki farklı dönem sonu notu alamaz.
Bir sınıftaki her öğrenciyi ebeveynleriyle eşleyen bağıntı ise fonksiyon
değildir: öğrencilerin sıfır, iki veya daha çok ebeveyni olabilir.
\end{ex}

\begin{explain}
Fonksiyonları, her olası argümandaki değerlerinin ne olduğunu kesin bir
biçimde belirterek tanımlayabiliriz. Bunun farklı yolları arasında bir formül
vermek, değeri hesaplayan bir yöntem betimlemek veya her argümandaki değerleri
listelemek bulunur. Fonksiyonları nasıl tanımlarsak tanımlayalım, her argüman
için bir ve yalnızca bir değer belirttiğimizden emin olmalıyız.
\end{explain}


\begin{ex}
$f(x)=x+1$ olacak biçimde bir $f \colon \Nat \to \Nat$ fonksiyonu
tanımlayalım. Bu tanım, $f$'yi doğal sayıları girdi olarak alıp doğal sayıları
çıktı olarak veren bir fonksiyon olarak belirler. Bir $x$ doğal sayısı
verildiğinde $f$'nin, onun ardılı olan $x+1$'i vereceğini söyler.
Bu durumda değer kümesi $\Nat$, $f$'nin görüntü kümesi değildir; çünkü $0$
doğal sayısı hiçbir doğal sayının ardılı değildir. $f$'nin görüntü kümesi,
bütün pozitif tam sayıların kümesi $\Int^{+}$'dır.
\end{ex}

\begin{ex}\ollabel{examplefunext}
$g(x)=x+2-1$ olacak biçimde bir $g \colon \Nat \to \Nat$ fonksiyonu
tanımlayalım. Bu, $g$'nin doğal sayıları girdi olarak alıp doğal sayıları
çıktı olarak veren bir fonksiyon olduğunu söyler. Bir $x$ doğal sayısı
verildiğinde $g$, $x$'in ardılının ardılının öncülünü, yani $x+1$'i verir.
\end{ex}

\begin{explain}
Az önce farklı \emph{tanımlara} sahip iki $f$ ve $g$ fonksiyonunu ele aldık.
Ancak bunlar \emph{aynı fonksiyondur}. Gerçekten her $n$ doğal sayısı için
$f(n)=n+1=n+2-1=g(n)$ olur. Başka bir deyişle $f$ ve $g$ tanımlarımız,
farklı denklemler aracılığıyla aynı eşlemeyi belirtir. Öyleyse örtük olarak
fonksiyonlara ilişkin bir kaplamsallık ilkesine dayanıyoruz:
\[
  \text{eğer }\forall x\, f(x) = g(x)\text{ ise }f = g
\]
Burada $f$ ile $g$'nin aynı tanım ve değer kümelerine sahip olduğu
varsayılır.
\end{explain}

\begin{ex}
Fonksiyonları durumlara göre de tanımlayabiliriz. Örneğin
$h \colon \Nat \to \Nat$ fonksiyonunu
\[
h(x) =
\begin{cases}
  \frac{x}{2} & \text{$x$ çiftse} \\
  \frac{x+1}{2} & \text{$x$ tekse.}
\end{cases}
\]
biçiminde tanımlayabiliriz. Her doğal sayı ya çift ya da tek olduğundan, bu
fonksiyonun çıktısı her zaman bir doğal sayı olacaktır. Bir fonksiyonu
durumlara göre tanımladığınızda her olası girdinin tam olarak bir duruma
girmesi gerektiğini unutmayın. Bazı durumlarda bunun için durumların bütün
olasılıkları tükettiğinin ve birbirini dışladığının ispatlanması gerekir.
\end{ex}

\end{document}
